\subsection{Особенности адресации (Буфер драйвера)}
Важно: \hyperlink{term:ledbar}{Ledbar} работает с \hyperlink{term:buffer}{буфером} \textbf{из 29 элементов}. Поэтому инициализация должна задавать размерность \verb|29|.
\begin{quote}
\textbf{Требование:} \verb|Ledbar.new(29)|.
\end{quote}

Типовой шаблон инициализации выглядит так:

\begin{codefile}{lua/examples/05_leds/leds_buffer_init.lua}
\end{codefile}

\textbf{Пояснение.} В примере создаётся объект \verb|Ledbar| на 29 светодиодов и функция \verb|changeColor(col)|, которая проходит по индексам \verb|0..28| и вызывает \verb|leds:set(i, ...)|. Основной принцип — один раз инициализировать буфер нужного размера и затем менять цвет сразу всей ленты через обход по индексам.
